\noindent These tables summarize the homogeneous sectors generated by \(\sigma_n\) and the way the sigma-conjugate, reverse, and combined operators act on each grade.
\begin{center}
\captionsetup{type=table,font=small,labelfont=bf,skip=4pt}
\captionof{table}{Sigma grades and representative basis elements.}\label{tab:sigma-grade-components-app}
\renewcommand{\arraystretch}{1.18}
\setlength{\tabcolsep}{4pt}
\small
\begin{tabularx}{\columnwidth}{@{}c l >{\raggedright\arraybackslash}X@{}}
\toprule
grade & sector & representative elements \\
\midrule
\(0\) & scalar & \(1\) \\
\(1\) & vector & \(\sigma_1,\ \sigma_2,\ \sigma_3\) \\
\(2\) & bivector & \(i\sigv=\{\sigma_2\sigma_3,\ \sigma_3\sigma_1,\ \sigma_1\sigma_2\}\) \\
\(3\) & pseudo-scalar & \(i=\sigma_1\sigma_2\sigma_3\) \\
\bottomrule
\end{tabularx}
\end{center}

\begin{center}
\captionsetup{type=table,font=small,labelfont=bf,skip=4pt}
\captionof{table}{Signs of the sigma-conjugate, reverse, and combined operators on homogeneous grades.}\label{tab:grade-signs}
\renewcommand{\arraystretch}{1.18}
\setlength{\tabcolsep}{4pt}
\small
\begin{tabularx}{\columnwidth}{@{}l >{\raggedright\arraybackslash}X >{\raggedright\arraybackslash}X >{\raggedright\arraybackslash}X@{}}
\toprule
Grade & sigma conjugate \((\cdot)^*\) & reverse operator \((\cdot)^{\mathrm R}\) & combined \((\cdot)^{*\mathrm R}\) \\
\midrule
scalar & \(+\) & \(+\) & \(+\) \\
vector & \(-\) & \(+\) & \(-\) \\
bivector & \(+\) & \(-\) & \(-\) \\
pseudo-scalar & \(-\) & \(-\) & \(+\) \\
\bottomrule
\end{tabularx}
\end{center}
